% Part: normal-modal-logic
% Chapter: filtrations
% Section: introduction

\documentclass[../../../include/open-logic-section]{subfiles}

\begin{document}

\olfileid{nml}{fil}{int}

\olsection{引言}

关于一个逻辑，始终重要的一个问题是它是否可判定，
即是否存在一个能行程序来回答“该公式是否有效”。
命题逻辑是可判定的：我们可以通过构造真值表来能行地测试一个公式是否是重言式，
并且对于给定的公式，真值表是有穷的。
但我们显然不能测试一个模态公式是否在所有模型中为真，
因为模型有无穷多个。
不过，我们可以列出与给定公式相关的所有有穷模型，
因为只有对公式中实际出现的命题变元指派世界的子集才是相关的。
如果可达关系固定，所有可能的不同赋值 $V(p)$ 就是 $W$ 的所有子集，
而如果 $\card{W} = n$，则这样的子集有 $2^n$ 个。
如果我们的公式 $!A$ 包含 $m$ 个命题变元，那么具有 $n$ 个世界的不同模型就有 $2^{nm}$ 个。
对其中每一个模型，我们只需通过在每个世界中计算 $!A$ 的真值，
就可以测试 $!A$ 是否在所有世界上为真。
当然，我们还必须检查所有可能的可达关系，
但 $n$ 个世界上也只有有穷多种关系（具体而言，是 $W \times W$ 的子集数，即 $2^{n^2}$）。

如果我们关心的不是逻辑 $\Log{K}$，而是由某个模型类（例如，自反传递模型）定义的逻辑，
我们还必须能测试可达关系是否属于所需的类型。
只要我们所关心的框架类可由模态公式定义（例如，通过测试公式 \Ax{T} 和 \Ax{4} 在该框架中是否有效），我们就能做到这一点。
因此，我们的想法会是：遍历所有有穷框架，
测试每个框架是否属于我们感兴趣的类；
然后列出该框架上所有可能的模型，
并测试 $!A$ 在其中是否为真。
如果不是，就停止：$!A$ 在我们感兴趣的模型类中不是有效的。

这个想法有一个问题：我们不知道何时（如果真有那一刻的话）可以停止搜索。
如果公式有一个有穷反模型，我们的过程会找到它。
但如果它没有有穷反模型，我们就得不到答案。
该公式可能是有效的（根本没有反模型），
或者它可能只有一个无穷的反模型，而我们永远不会检查到。
如果我们能证明，每个有反模型的公式都有一个有穷反模型，那么这个问题就能得到解决。
如果情况如此，我们就说该逻辑具有\emph{有穷模型性质}。

但我们如何证明一个逻辑具有有穷模型性质呢？
一种方法是找到一种将 $!A$ 的无穷（反）模型转化为有穷模型的方法。
如果能做到这一点，那么每当存在一个 $!A$ 在其中不为真的模型时，
转化得到的有穷模型也会使 $!A$ 不为真。
那个有穷模型将会出现在我们所有有穷模型的列表中，
而我们最终将对每个不是有效的公式判定它不是有效的。
如果公式是有效的，我们的过程将不会终止。
如果我们还能进一步证明，我们过程所给出的有穷模型存在一个最大大小，
并且这个最大大小只依赖于公式 $!A$，
那么我们在搜索反模型时，只需测试不超过该大小的有穷模型。
如果到那时还没有找到反模型，那就说明没有反模型。
这样一来，我们的过程实际上就能对任意公式 $!A$ 判定“$!A$ 是否有效？”这个问题。

一种常能成功地将无穷结构转化为有穷结构的策略是，将结构中在相关方面行为相同的元素“等同”起来。
如果 $\mModel{M}$ 中有无穷多个世界在相关方面行为相同，
那么我们或许就能将所有世界收集到有穷多个（可能无穷的）此类世界的“类”中。
换句话说，我们应该以正确的方式划分世界集合，
即，使得每个划分块包含无穷多个世界，但只有有穷多个划分块。
然后我们定义一个新模型 $\mModel{M^*}$，其中的世界就是这些划分块。
旧模型中的有穷多个划分块给了我们新模型中的有穷多个世界，即一个有穷模型。
让我们称世界 $w$ 所在的划分块为 $[w]$。
我们希望 $\mSat{M}{!A}[w]$ 当且仅当
$\mSat{M^*}{!A}[{[w]}]$，因为我们希望如果旧模型是 $!A$ 的反模型，那么新模型也是。
这就要求我们以正确的方式来定义划分以及 $\mModel{M^*}$ 的可达关系。

为了看清具体怎么做，首先想象我们没有可达关系。$\mSat{M}{\Box !B}[w]$ 当且仅当存在某个 $v \in W$ 使得 $\mSat{M}{\Box !B}[v]$，对于 $\mModel{M^*}$ 也是如此，只是要把 $w$ 和 $v$ 换成 $[w]$ 和 $[v]$。作为初步的想法，如果两个世界 $u$ 和 $v$ 在 $\mModel{M}$ 中对所有命题变元的赋值一致，我们就说它们是等价的（属于同一个划分），即 $\mSat{M}{p}[u]$ 当且仅当 $\mSat{M}{p}[v]$。令 $V^*(p) = \Setabs{[w]}{\mSat{M}{p}[w]}$。我们的目标是证明 $\mSat{M}{!A}[w]$ 当且仅当 $\mSat{M^*}{!A}[{[w]}]$。显然，我们可以用归纳法来证明这一点：基始情形是 $!A \equiv p$。先假设 $\mSat{M}{p}[w]$。那么根据定义，$[w] \in V^*$，所以 $\mSat{M^*}{p}[{[w]}]$。现在假设 $\mSat{M^*}{p}[{[w]}]$。这意味着 $[w] \in V^*(p)$，即存在某个与 $w$ 等价的 $v$，使得 $\mSat{M}{p}[v]$。但“$w$ 等价于 $v$”意味着“$w$ 和 $v$ 使相同的命题变元为真”，所以 $\mSat{M}{p}[w]$。现在来看归纳步骤，例如 $!A \ident \lnot !B$ 的情形。那么 $\mSat{M}{\lnot !B}[w]$ 当且仅当 $\mSat/{M}{!B}[w]$ 当且仅当 $\mSat/{M^*}{!B}[{[w]}]$（由归纳假设）当且仅当 $\mSat{M^*}{\lnot !B}[{[w]}]$。其他非模态算子的情形类似。对于 $\Box$ 也成立：假设 $\mSat{M^*}{\Box
  !B}[{[w]}]$。这意味着对于每个 $[u]$，$\mSat{M^*}{!B}[{[u]}]$。由归纳假设，对于每个 $u$，$\mSat{M}{!B}[u]$。因此，$\mSat{M}{\Box !B}[w]$。

在一般情形下，即我们还需要为 $\mModel{M^*}$ 定义可达关系时，事情就变得更复杂了。如果模型 $\mModel{M^*}$ 的可达关系 $R^*$ 满足使上述归纳证明成立所需的条件，我们就称该模型为一个\emph{滤过}。那么，任何滤过 $\mModel{M^*}$ 都会使得 $!A$ 在 $[w]$ 上为真，当且仅当 $\mModel{M}$ 使得 $!A$ 在 $w$ 上为真。然而，现在我们还必须证明滤过\emph{存在}，也就是说，我们可以定义 $R^*$ 使其满足所需的条件。但是，为了使这能行得通，我们必须要求世界 $u$、$v$ 被视为等价的条件，不仅仅是它们对所有命题变元的赋值一致，而且是对 $!A$ 的所有子公式的赋值也一致。由于 $!A$ 只有有穷多个子公式，这仍然能保证滤过是有穷的。定义滤过的方法不止一种，而为了确保滤过的可达关系满足所需的性质（例如自反、传递等），我们必须在 $R^*$ 的定义上发挥创造性。

\end{document}